\section{语言模型}

\subsection{什么是语言模型}

\begin{importantbox}{语言模型的定义}
\textbf{语言模型}（Language Model, LM）是一个能够对文本序列中下一个词的出现概率进行预测的系统。形式化地，给定前文 $x^{(1)}, x^{(2)}, \ldots, x^{(t)}$，语言模型输出下一个词的条件概率分布：
$$P(x^{(t+1)} \mid x^{(1)}, x^{(2)}, \ldots, x^{(t)})$$
该概率分布满足 $\sum_{w \in V} P(w \mid x^{(1)}, \ldots, x^{(t)}) = 1$，其中 $V$ 是词汇表。
\end{importantbox}

等价地，语言模型可以为任意一段文本 $x^{(1)}, \ldots, x^{(T)}$ 赋予一个概率。通过链式法则：

$$P(x^{(1)}, \ldots, x^{(T)}) = P(x^{(1)}) \cdot P(x^{(2)} \mid x^{(1)}) \cdot P(x^{(3)} \mid x^{(1)}, x^{(2)}) \cdots P(x^{(T)} \mid x^{(1)}, \ldots, x^{(T-1)})$$

$$= \prod_{t=1}^{T} P(x^{(t)} \mid x^{(1)}, \ldots, x^{(t-1)})$$

\begin{figure}[H]
\centering
\includegraphics[width=0.95\textwidth]{slides-images/slide-013.jpg}
\caption{语言模型的基本任务：给定前文``the students opened their''，预测下一个词的概率分布\protect\footnotemark}
\end{figure}
\footnotetext{Slides 第14页。}

\subsection{语言模型无处不在}

Manning 强调，语言模型并非 ChatGPT 时代的发明。它的历史可以追溯到 1950 年代 Claude Shannon 的信息论工作，至少从 1980 年代起就已经是 NLP 的核心技术。

\begin{figure}[H]
\centering
\includegraphics[width=0.95\textwidth]{slides-images/slide-015.jpg}
\caption{日常生活中的语言模型：手机输入法的下一词预测、Google 搜索的自动补全都是语言模型的应用\protect\footnotemark}
\end{figure}
\footnotetext{Slides 第16页。}

\subsection{N-gram 语言模型}

在神经网络语言模型出现之前（约 1975--2012 年），\textbf{N-gram 语言模型}是主流方法。

\subsubsection{基本思想}

N-gram 模型做了一个\textbf{马尔可夫假设}（Markov Assumption）：预测下一个词时，只使用前 $n-1$ 个词作为上下文，丢弃更早的信息。

$$P(x^{(t+1)} \mid x^{(1)}, \ldots, x^{(t)}) \approx P(x^{(t+1)} \mid x^{(t-n+2)}, \ldots, x^{(t)})$$

概率估计通过计数得到：

$$P(w \mid \text{students opened their}) = \frac{\text{count}(\text{students opened their } w)}{\text{count}(\text{students opened their})}$$

\begin{figure}[H]
\centering
\includegraphics[width=0.95\textwidth]{slides-images/slide-017.jpg}
\caption{4-gram 语言模型示例：基于``students opened their''预测下一个词，books 的概率为 0.4，exams 为 0.1\protect\footnotemark}
\end{figure}
\footnotetext{Slides 第18页。}

\subsubsection{N-gram 模型的问题}

\begin{warningbox}{N-gram 模型的两大致命缺陷}
\begin{enumerate}
    \item \textbf{稀疏性问题}（Sparsity）：很多合理的 n-gram 组合在训练语料中从未出现过，导致概率估计为零。而概率为零是灾难性的——任何涉及该 n-gram 的计算都会变为零。解决方法包括加法平滑（Laplace Smoothing）和回退（Backoff）策略。
    \item \textbf{存储问题}（Storage）：需要存储语料中所有观察到的 n-gram 的计数。增大 $n$ 或增大语料规模都会使存储需求指数级增长。实践中 $n$ 通常不超过 5。
\end{enumerate}
\end{warningbox}

\begin{figure}[H]
\centering
\includegraphics[width=0.95\textwidth]{slides-images/slide-020.jpg}
\caption{N-gram 模型的存储问题：增大 $n$ 或增大语料都会增加模型大小\protect\footnotemark}
\end{figure}
\footnotetext{Slides 第21页。}

\subsubsection{用 N-gram 模型生成文本}

尽管 N-gram 模型很``粗糙''，但它已经能生成看起来有点像样的文本。Manning 展示了用 trigram 模型在 200 万词语料上生成的文本：

\begin{figure}[H]
\centering
\includegraphics[width=0.95\textwidth]{slides-images/slide-025.jpg}
\caption{Trigram 语言模型生成的文本：语法大致正确（surprisingly grammatical），但语义不连贯（incoherent）\protect\footnotemark}
\end{figure}
\footnotetext{Slides 第26页。}

\begin{knowledgebox}{Scale 的故事古已有之}
Manning 有一个精彩的观察：今天人们说的``更多数据、更大模型就能解决一切''，与 2010 年代初人们对 N-gram 模型的态度如出一辙——当时人们也在说，只要有更大的语料和更高阶的 N-gram，性能就会持续提升。但事实证明，\textbf{更好的模型架构}同样重要。
\end{knowledgebox}

\subsection{固定窗口神经语言模型}

N-gram 模型的自然演进是使用神经网络来替代简单的计数。Yoshua Bengio 在 2003 年提出了\textbf{固定窗口神经语言模型}（Fixed-Window Neural Language Model）。

\begin{figure}[H]
\centering
\includegraphics[width=0.95\textwidth]{slides-images/slide-028.jpg}
\caption{固定窗口神经语言模型：将前 $n$ 个词的词嵌入拼接，通过隐藏层和 softmax 预测下一个词\protect\footnotemark}
\end{figure}
\footnotetext{Slides 第29页。}

该模型的工作流程如下：

\begin{enumerate}
    \item 将窗口内的每个词通过词嵌入矩阵 $E$ 转换为向量 $e^{(i)}$
    \item 将所有词嵌入拼接为一个长向量
    \item 通过隐藏层：$h = f(W \cdot [e^{(1)}; e^{(2)}; \ldots; e^{(n)}] + b_1)$
    \item 通过输出层和 softmax：$\hat{y} = \text{softmax}(U \cdot h + b_2)$
\end{enumerate}

\begin{importantbox}{固定窗口模型的优势与局限}
\textbf{优势}：
\begin{itemize}
    \item 无稀疏性问题——不需要存储所有 n-gram 计数
    \item 存储成本更低——只需存储网络参数
    \item 通过词嵌入实现了一定的泛化能力
\end{itemize}
\textbf{局限}：
\begin{itemize}
    \item 仍然使用固定窗口，仍有马尔可夫假设的限制
    \item 不同位置的词使用权重矩阵 $W$ 的不同部分处理，没有参数共享——同一个词出现在不同位置时，学到的知识无法迁移
    \item 窗口大小的选择是个难题——没有任何固定窗口是``足够大''的
\end{itemize}
\end{importantbox}

\subsection{本章小结}

\begin{itemize}
    \item 语言模型是预测下一个词的概率分布的系统，是 NLP 的核心技术
    \item N-gram 语言模型通过计数和马尔可夫假设工作，简单但受限于稀疏性和存储问题
    \item 固定窗口神经语言模型解决了稀疏性问题，但仍受限于固定上下文长度和缺乏参数共享
    \item 我们需要一种能处理\textbf{任意长度}输入且\textbf{共享参数}的新架构——这就引出了循环神经网络
\end{itemize}

%% ========== Part 3: 循环神经网络 ==========

